\section*{Bab \ref{ch:g10:why-clouds-form} --- Mengapa Awan Terbentuk}

\begin{solution}{exo:g10:why-clouds-form:1}
\omterm{def:g10:why-clouds-form:humidity}{Kelembapan relatif} adalah kandungan uap air udara dibagi kandungan \omterm{def:g10:why-clouds-form:humidity}{jenuhnya}
pada suhu dan tekanan yang sama, dinyatakan dalam persen. \omterm{def:g10:why-clouds-form:dew-point}{Titik embun} adalah
suhu yang harus dicapai udara bila didinginkan, dengan kandungan air dan
tekanannya tetap, agar menjadi \omterm{def:g10:why-clouds-form:humidity}{jenuh}.
\end{solution}

\begin{solution}{exo:g10:why-clouds-form:2}
\num{3.8}, \num{7.6}, \num{14.7}, dan \qty{27.2}{g/kg}. Setiap
\qty{10}{\degreeCelsius} mengalikannya dengan sekitar 2 ($7.6/3.8 = 2.0$,
$14.7/7.6 \approx 1.9$, $27.2/14.7 \approx 1.9$).
\end{solution}

\begin{solution}{exo:g10:why-clouds-form:3}
Kandungan \omterm{def:g10:why-clouds-form:humidity}{jenuh} pada \qty{24}{\degreeCelsius}: \qty{18.9}{g/kg}; \omterm{def:g10:why-clouds-form:humidity}{kelembapan
relatif} $10.0/18.9 \approx \qty{53}{\percent}$. Kandungan \omterm{def:g10:why-clouds-form:humidity}{jenuhnya}
\qty{10.0}{g/kg} pada \qty{14}{\degreeCelsius}: itulah \omterm{def:g10:why-clouds-form:dew-point}{titik embunnya}.
\end{solution}

\begin{solution}{exo:g10:why-clouds-form:4}
Pada musim panas udara lembap dan \omterm{def:g10:why-clouds-form:dew-point}{titik embunnya} jauh di atas suhu gelas:
udara yang menyentuh gelas didinginkan sampai di bawah \omterm{def:g10:why-clouds-form:dew-point}{titik embunnya} dan air
mengembun pada gelas. Udara musim dingin yang sangat kering di dalam ruangan
memiliki \omterm{def:g10:why-clouds-form:dew-point}{titik embun} yang lebih rendah daripada suhu gelas: meskipun
didinginkan oleh gelas, udara itu tetap tidak \omterm{def:g10:why-clouds-form:humidity}{jenuh}.
\end{solution}

\begin{solution}{exo:g10:why-clouds-form:5}
Pemanasan dari bawah (kumulus sore hari di atas tanah yang tersinari
Matahari); pegunungan (angin dipaksa naik melewati punggung pegunungan); front
(udara hangat naik di atas udara dingin); konvergensi (angin-angin mengalir
menyatu di dekat khatulistiwa atau ke dalam depresi).
\end{solution}

\begin{solution}{exo:g10:why-clouds-form:6}
$125 \times (30 - 18) = \qty{1500}{m}$; $125 \times (22 - 20) =
\qty{250}{m}$. Hari kedua: awan terbentuk hanya \qty{250}{m} di atas
tanah.
\end{solution}

\begin{solution}{exo:g10:why-clouds-form:7}
Pada \qty{1}{km}: $18 - 9.8 = \qty{8.2}{\degreeCelsius}$; pada \qty{2}{km}:
$18 - 19.6 = \qty{-1.6}{\degreeCelsius}$. Ketika turun, parsel itu menghangat
dengan laju yang sama dan mencapai tanah kembali pada \qty{18}{\degreeCelsius}.
\end{solution}

\begin{solution}{exo:g10:why-clouds-form:8}
Dengan \qty{4}{\degreeCelsius/km}, baik parsel tak \omterm{def:g10:why-clouds-form:humidity}{jenuh}
(\qty{9.8}{\degreeCelsius/km}) maupun parsel \omterm{def:g10:why-clouds-form:humidity}{jenuh}
(\qty{5}{\degreeCelsius/km}) mendingin lebih cepat daripada sekitarnya: \omterm{def:g10:why-clouds-form:stability}{stabil}
dalam kedua kasus. Dengan \qty{12}{\degreeCelsius/km}, gelembung tak \omterm{def:g10:why-clouds-form:humidity}{jenuh} pun
mendingin lebih lambat daripada sekitarnya: gelembung itu tetap lebih hangat
dan terus naik, dan gelembung-gelembung udara hangat naik dengan bebas.
\end{solution}

\begin{solution}{exo:g10:why-clouds-form:9}
$14.7 - 7.6 = \qty{7.1}{g}$ per kilogram udara. \omterm{def:g10:earth-energy-budget:latent-heat}{Kalor laten} itu menghangatkan
udara tempat air itu mengembun, sehingga memperlambat pendinginannya.
\end{solution}

\begin{solution}{exo:g10:why-clouds-form:10}
Volume bertambah sebanding dengan pangkat tiga diameter:
$(0.5/0.02)^3 = 25^3 = 15\,625$ titik awan untuk satu tetes gerimis;
$(2/0.02)^3 = 100^3 = 1\,000\,000$ untuk satu tetes hujan.
\end{solution}

\begin{solution}{exo:g10:why-clouds-form:11}
Di bawah langit cerah, tanah memancarkan inframerah dan mendingin, begitu pula
udara yang menyentuhnya. Udara dingin ini lebih rapat daripada udara di
atasnya; di lereng-lereng, udara itu mengalir turun dan berkumpul di dasar
lembah. Di sana udara itu terus mendingin sepanjang malam sampai mencapai
\omterm{def:g10:why-clouds-form:dew-point}{titik embunnya}, dan titik-titik air terbentuk di dalam udara itu sendiri:
kabut. Di puncak bukit udara tidak sedingin itu dan tetap di atas \omterm{def:g10:why-clouds-form:dew-point}{titik
embunnya}.
\end{solution}

\begin{solution}{exo:g10:why-clouds-form:12}
Pada ketinggian $z$ suhunya $T - a z$ dan \omterm{def:g10:why-clouds-form:dew-point}{titik embunnya} $T_d - b z$;
keduanya sama bila $(a - b) z = T - T_d$, jadi $z = (T - T_d)/(a - b)$. Dengan
$b = 1.7$: $a - b = 8.1$, dan $1000/8.1 \approx \qty{123}{m}$ per derajat.
\end{solution}

\begin{solution}{exo:g10:why-clouds-form:13}
\omterm{def:g10:why-clouds-form:condensation-level}{Aras kondensasi}: $125 \times (15 - 11) = \qty{500}{m}$; suhu di sana:
$15 - 9.8 \times 0.5 \approx \qty{10.1}{\degreeCelsius}$. Puncak: $10.1 - 5
\times 1.5 = \qty{2.6}{\degreeCelsius}$. Turun \qty{2}{km} dengan laju kering:
$2.6 + 9.8 \times 2 \approx \qty{22.2}{\degreeCelsius}$, sekitar tujuh derajat
lebih hangat daripada di sisi hadapan angin.
\end{solution}

\begin{solution}{exo:g10:why-clouds-form:14}
Pada \qty{-20}{\degreeCelsius}: $126/103 \approx 1.22$, lewat \omterm{def:g10:why-clouds-form:humidity}{jenuh}
\qty{22}{\percent} bagi es; pada \qty{-10}{\degreeCelsius}: $287/260 \approx
1.10$, \qty{10}{\percent}. Di antara suhu-suhu ini, makin dingin awan, makin
besar keunggulan es: proses itu bekerja paling baik di bagian atas awan yang
dingin, tempat titik-titik air dan kristal es bercampur.
\end{solution}

\begin{solution}{exo:g10:why-clouds-form:15}
Tak \omterm{def:g10:why-clouds-form:humidity}{jenuh}: parsel mendingin \qty{9.8}{\degreeCelsius}, sekitarnya
\qty{6.5}{\degreeCelsius}: parsel \qty{3.3}{\degreeCelsius} lebih dingin dan
turun kembali. \omterm{def:g10:why-clouds-form:humidity}{Jenuh}: parsel mendingin \qty{5}{\degreeCelsius}, jadi
\qty{1.5}{\degreeCelsius} lebih hangat dan terus naik. Udara itu \omterm{def:g10:why-clouds-form:stability}{tidak stabil}
hanya dengan satu syarat: udara yang naik itu \omterm{def:g10:why-clouds-form:humidity}{jenuh}.
\end{solution}

\begin{solution}{pb:g10:why-clouds-form:1}
\textbf{1.} \qty{14.7}{g/kg}.

\textbf{2.} $11.4/14.7 \approx \qty{78}{\percent}$.

\textbf{3.} \qty{16}{\degreeCelsius}, tempat kandungan \omterm{def:g10:why-clouds-form:humidity}{jenuhnya}
\qty{11.4}{g/kg}.

\textbf{4.} $11.4/24.1 \approx \qty{47}{\percent}$.

\textbf{5.} Tidak: udara itu tetap mengandung \qty{11.4}{g/kg}. \omterm{def:g10:why-clouds-form:humidity}{Kelembapan
relatifnya} turun karena, setelah lebih hangat, udara itu dapat menampung lebih
banyak.

\textbf{6.} Pada \omterm{def:g10:why-clouds-form:dew-point}{titik embunnya}, \qty{16}{\degreeCelsius}.

\textbf{7.} $11.4 - 8.7 = \qty{2.7}{g}$ per kilogram udara.

\textbf{8.} Di bawah langit malam yang cerah, rumput dan atap memancarkan
inframerah dan menjadi lebih dingin daripada udara; udara yang menyentuhnya
didinginkan sampai di bawah \omterm{def:g10:why-clouds-form:dew-point}{titik embunnya} sebelum udara lainnya.

\textbf{9.} $(28.0 - 13.4)/1.5 \approx \qty{9.7}{\degreeCelsius/km}$,
praktis sama dengan laju kering: gelembung-gelembung sore hari telah mengaduk
lapisan ini.

\textbf{10.} $(13.4 - 0.4)/2.0 = \qty{6.5}{\degreeCelsius/km}$.

\textbf{11.} \omterm{def:g10:why-clouds-form:stability}{Stabil} bagi parsel tak \omterm{def:g10:why-clouds-form:humidity}{jenuh}, yang mendingin lebih cepat
(\qty{9.8}{\degreeCelsius/km}); \omterm{def:g10:why-clouds-form:stability}{tidak stabil} bagi parsel \omterm{def:g10:why-clouds-form:humidity}{jenuh}, yang mendingin
lebih lambat (\qty{5}{\degreeCelsius/km}).

\textbf{12.} Suhu naik terhadap ketinggian, dari \num{-0.3} menjadi
\qty{1.2}{\degreeCelsius}: lapisan yang sangat \omterm{def:g10:why-clouds-form:stability}{stabil}, yang bekerja sebagai
tutup. Awan yang tumbuh ke atas berhenti di sana: puncaknya berada di dekat
\qty{4}{km}.

\textbf{13.} $28 - 16 = \qty{12}{\degreeCelsius}$.

\textbf{14.} Suhu $28 - 9.8 = \qty{18.2}{\degreeCelsius}$; \omterm{def:g10:why-clouds-form:dew-point}{titik embun}
$16 - 1.8 = \qty{14.2}{\degreeCelsius}$.

\textbf{15.} Dari 12 menjadi $18.2 - 14.2 = 4.0$: selisihnya menyusut
\qty{8}{\degreeCelsius}, dengan laju \qty{8}{\degreeCelsius} per kilometer.

\textbf{16.} Selisih satu derajat tertutup dalam $1/8$ km, yaitu
\qty{125}{m}.

\textbf{17.} $125 \times (20 - 16) = \qty{500}{m}$.

\textbf{18.} \omterm{def:g10:why-clouds-form:dew-point}{Titik embun} tetap \qty{16}{\degreeCelsius}, sedangkan udara di
dekat tanah menghangat: selisihnya bertambah, begitu pula ketinggian tempat
selisih itu tertutup.

\textbf{19.} $28 - 9.8\,z = 16 - 1.8\,z$.

\textbf{20.} $8\,z = 12$, jadi $z = 1.5$: dasar kumulus sore hari berada
\textbf{sekitar \qty{1.5}{km}} di atas dataran (perhitungan yang tepat dengan
sifat-sifat air memberikan \qty{1.51}{km}).
\end{solution}
